% --- 01 Motivation: quantum threat -------------------------------------------
\begin{frame}{Η κβαντική απειλή}
  \begin{columns}[T,onlytextwidth]
    \begin{column}{0.56\textwidth}
      \begin{itemize}
        \item Ο αλγόριθμος του \en{Shor} σπάει \en{RSA} / \en{ECC} σε πολυωνυμικό χρόνο~\cite{Shor1997}
        \item \HL{“\en{Harvest now, decrypt later}”}: τα δεδομένα κλέβονται σήμερα, αποκρυπτογραφούνται αύριο
        \item Εντολή \en{NSA CNSA 2.0}: μετάβαση σε \en{PQC} έως το \HL{2035}~\cite{NSA-CNSA2}
      \end{itemize}
      \vspace{0.6em}
      \pause
      \begin{block}{Γιατί \HQC;}
        Τον Μάρτιο 2025 το \en{NIST} επέλεξε το \HQC ως πρότυπο \en{KEM}, εναλλακτικό στο \en{ML-KEM}, βασισμένο σε κωδικοθεωρία~\cite{NIST-HQC-selection2025}.
      \end{block}
    \end{column}
    \begin{column}{0.44\textwidth}
      \begin{exampleblock}{Διαφοροποίηση κινδύνου}
        Το \HQC στηρίζεται σε \emph{διαφορετική} υπόθεση σκληρότητας (αποκωδικοποίηση συγκλινούσας) απ' ό,τι τα πλεγματικά, αντιστάθμιση σε μια μελλοντική κρυπτανάλυση.
      \end{exampleblock}
    \end{column}
  \end{columns}
\end{frame}
